% Ad Familiares 3.11 --- headnote
% ad-familiares-03-11, June 50 BC, camp by the river Pyramus in Cilicia

\textit{Cicero to Appius Claudius Pulcher, written from
camp by the river Pyramus in Cilicia in late June 50 BC
(Perseus dateline: \textit{Scr. in castris ad Pyramum
post Id. Iun. 704 (50)}). The letter answers two of
Appius's at once, brought up from Tarsus by Q.
Servilius: the earlier dated to the Nones of April,
the later undated. Its principal occasion is the news
of Appius's acquittal on the \textit{maiestas} charge
that Dolabella had pressed against him in Rome. The
prosecution was a piece of political theatre that
Cicero suppresses an awkwardness about: Dolabella was
at the same moment Tullia's prospective husband, the
engagement having been concluded in Cicero's absence
by his wife and daughter. The letter passes over
this entirely. It is congratulatory, careful, and just
a little stylised --- the salute already addresses
Appius as ``censor, I hope,'' anticipating the
elevation that the acquittal had cleared the way for.}

\textit{Section 2 contains a small and characteristic
piece of Ciceronian forensic comment: \textit{maiestas},
Cicero observes, is Sulla's vague catch-all, where
declamation against anyone is hard to bound; but
\textit{ambitus} --- the alternative charge --- is so
plain a matter of fact that it must either be brought
shamefully or defended. He pictures himself in court
beside Appius, drawing peals of laughter against the
prosecution he was not present to deride. The praise
of Pompey and Brutus in section 3 is set carefully: of
all kinsmen and friends, ``the one the first man of
every age and nation, the other long since the first
of our young men and soon, as I hope, of our state.''
The shadow of the gathering crisis is already on the
sentence.}

\textit{Section 4 turns to Appius's second letter,
which had sketched the political weather at Rome.
Cicero reads it with relief: the dangers are lighter,
the resources greater --- ``if, as you write, all the
forces of the state have ranged themselves under
Pompey's command.'' Section 5 closes with the
private byplay of the correspondence: a complaint
that one of Cicero's earlier letters had been called
``rather peevish'' and not ``eloquent,'' answered with
the joke that, as Aristarchus refuses to credit Homer
with any line he does not approve, so anything
unpolished must not be Cicero's. The farewell looks
forward to the censorship Cicero hopes Appius will
hold; in fact the censors of 50 BC would be Appius
himself and L. Calpurnius Piso.}
